% ============================================================
% Paper 92 (DE): Donaldsons Theorie der 4-Mannigfaltigkeiten
%                und offene Probleme in der Eichtheorie
% Autor: Michael Fuhrmann
% Projekt: Specialist Mathematics Project
% Build: 169
% Datum: 2026-03-12
% ============================================================
\documentclass[12pt,a4paper]{amsart}

\usepackage{amsmath,amssymb,amsthm,mathtools}
\usepackage{enumitem}
\usepackage{booktabs}
\usepackage{array}
\usepackage[hidelinks]{hyperref}
\usepackage[utf8]{inputenc}
\usepackage[T1]{fontenc}
\usepackage[ngerman]{babel}
\usepackage{geometry}
\geometry{margin=2.5cm}

% ---- Theorem-Umgebungen -------------------------------------
\newtheorem{theorem}{Theorem}[section]
\newtheorem{lemma}[theorem]{Lemma}
\newtheorem{corollary}[theorem]{Korollar}
\newtheorem{proposition}[theorem]{Proposition}
\theoremstyle{definition}
\newtheorem{definition}[theorem]{Definition}
\newtheorem{example}[theorem]{Beispiel}
\theoremstyle{remark}
\newtheorem{remark}[theorem]{Bemerkung}
\newtheorem{conjecture}[theorem]{Vermutung}
% Deutsche Aliasse
\newtheorem{satz}[theorem]{Satz}
\newtheorem{korollar}[theorem]{Korollar}
\newtheorem{vermutung}[theorem]{Vermutung}
\newtheorem{bemerkung}[theorem]{Bemerkung}
\newtheorem{definition_de}[theorem]{Definition}

% ---- Operatoren ---------------------------------------------
\DeclareMathOperator{\rank}{rank}
\DeclareMathOperator{\sign}{sign}
\DeclareMathOperator{\Tr}{Tr}
\DeclareMathOperator{\SW}{SW}
\newcommand{\CP}{\mathbb{CP}}
\newcommand{\ZZ}{\mathbb{Z}}
\newcommand{\RR}{\mathbb{R}}
\newcommand{\CC}{\mathbb{C}}

% ---- Zusammenfassung ----------------------------------------
\begin{document}
\begin{abstract}
Wir geben einen rigoros mathematischen Überblick über Donaldsons
grundlegende Resultate zur glatten Topologie 4-dimensionaler Mannigfaltig\-keiten.
Im Mittelpunkt steht der Satz von 1983, wonach die Schnittform einer
glatten, geschlossenen, einfach zusammenhängenden 4-Mannigfaltigkeit mit
definiter Schnittform über $\ZZ$ diagonalisierbar sein muss.
Der Beweis beruht auf dem Yang--Mills-Instanton-Modulraum und
Uhlenbecks Kompaktheits\-satz.
Wir besprechen ferner die Seiberg--Witten-Theorie (1994), exotische
$\RR^4$-Strukturen, Kirby-Kalkül, Freedmans topologische Klassifikation
sowie die bis heute offene glatte vierdimensionale Poincaré-Vermutung.
Wittens Vermutung---dass Seiberg--Witten-Invarianten die
Donaldson-Invarianten bestimmen---wird in der von Taubes bewiesenen
Form diskutiert.
\end{abstract}

\title{Donaldsons Theorie der 4-Mannigfaltigkeiten\\
\large und offene Probleme in der Eichtheorie}

\author{Michael Fuhrmann}
\address{Specialist Mathematics Project, Build 169}
\email{specialist-maths@project.local}
\date{2026-03-12}

\maketitle
\tableofcontents

% =============================================================
\section{Einleitung}
% =============================================================

Die Topologie glatter 4-Mannigfaltigkeiten ist dramatisch reicher---und
rätselhafter---als in jeder anderen Dimension.  In allen anderen Dimensionen
$n \neq 4$ sind die glatte und die topologische Kategorie gut verstanden:
Sie stimmen für $n \le 3$ überein und unterscheiden sich für $n \ge 5$ in
kontrollierter Weise durch Chirurgietheorie.  Dimension vier steht für sich.

Das Schlüsselereignis war Donaldsons Arbeit von 1983 \cite{Donaldson1983},
die zeigte, dass die Schnittform einer glatten, definiten, einfach
zusammenhängenden 4-Mannigfaltigkeit \emph{diagonalisierbar} sein muss.
Kombiniert mit Freedmans gleichzeitiger topologischer Klassifikation
\cite{Freedman1982} folgte die Existenz \emph{exotischer $\RR^4$}---glatter
Mannigfaltigkeiten, die homöomorph, aber nicht diffeomorph zu $\RR^4$ sind.

Die Werkzeuge stammen aus der Physik: Yang--Mills-Eichtheorie, anti-selbst-duale
Zusammenhänge (Instantonen) und Modulräume.  Ein Jahrzehnt später führten
Seiberg und Witten \cite{SeibergWitten1994} ein weit einfacheres
Gleichungssystem ein, das im Wesentlichen dieselben glatten Invarianten erfasst.

% =============================================================
\section{Schnittformen}
% =============================================================

Sei $X$ eine glatte, geschlossene, orientierte 4-Mannigfaltigkeit.
Die \emph{Schnittform} ist die symmetrische Bilinearform
\[
  Q_X \colon H_2(X;\ZZ) \times H_2(X;\ZZ) \longrightarrow \ZZ,
\]
definiert durch das Zählen orientierter Schnittpunkte transversal gewählter
Repräsentanten, oder äquivalent als Cup-Produkt-Paarung via Poincaré-Dualität.

\begin{definition_de}[Definite/Indefinite Form]
$Q_X$ heißt \emph{definit}, falls sie über $\RR$ positiv oder negativ definit
ist.  Andernfalls heißt sie \emph{indefinit}.  Rang $b_2 = \rank Q_X$,
Signatur $\sigma = b_2^+ - b_2^-$ und Parität klassifizieren $Q_X$ über $\RR$.
\end{definition_de}

\begin{example}
$\CP^2$ hat $Q = (1)$; $\overline{\CP}^2$ hat $Q = (-1)$;
$S^2 \times S^2$ hat $Q = \begin{pmatrix}0&1\\1&0\end{pmatrix}$ (hyperbolisch);
die $K3$-Fläche hat $Q = -2E_8 \oplus 3H$.
\end{example}

Definite Formen über $\ZZ$ umfassen die diagonale Form $\pm\mathbf{1}_n$,
aber auch $E_8$, $E_6$, $D_4$ und viele weitere nicht-diagonalisierbare
Formen bei großem Rang.

% =============================================================
\section{Freedmans topologische Klassifikation}
% =============================================================

\begin{satz}[Freedman, 1982 \cite{Freedman1982}]
Jede symmetrische unimodulare Bilinearform über $\ZZ$ wird als Schnittform
einer einfach zusammenhängenden geschlossenen topologischen 4-Mannigfaltigkeit
realisiert.  Ferner sind zwei solche Mannigfaltigkeiten genau dann
homöomorph, wenn sie dieselbe Schnittform und denselben
Kirby--Siebenmann-Invarianten $\mathrm{ks} \in \ZZ/2\ZZ$ haben.
\end{satz}

Insbesondere wird jede definite Form---darunter $E_8 \oplus E_8$,
$-E_8$ usw.---topologisch realisiert.  Die Frage lautet: Welche lassen sich
\emph{glatt} realisieren?

% =============================================================
\section{Yang--Mills-Eichtheorie}
% =============================================================

\subsection{Aufstellung}

Sei $E \to X$ ein hermitesches Vektorbündel vom Rang $r$ über einer
Riemannschen 4-Mannigfaltigkeit $(X,g)$ mit Strukturgruppe
$G = \mathrm{SU}(r)$.  Sei $\mathcal{A}$ der Raum der $G$-Zusammenhänge
auf $E$, und $\mathcal{G} = \Gamma(\Ad E)$ die Eichgruppe.

\begin{definition_de}[Yang--Mills-Funktional]
Das Yang--Mills-Funktional ist
\[
  \mathcal{YM}(A) = \int_X |F_A|^2\,d\mathrm{vol}_g,
\]
wobei $F_A \in \Omega^2(X, \ad E)$ die Krümmungs-2-Form von $A$ ist.
\end{definition_de}

\subsection{Selbstdualitätszerlegung}

In Dimension 4 erfüllt der Hodge-Stern $\star \colon \Omega^2 \to \Omega^2$
die Identität $\star^2 = \mathrm{id}$, was die orthogonale Zerlegung
\[
  \Omega^2(X) = \Omega^2_+(X) \oplus \Omega^2_-(X)
\]
liefert.  Mit $F_A = F_A^+ + F_A^-$ folgt durch Chern--Weil-Theorie:
\[
  \mathcal{YM}(A) = \|F_A^+\|^2 + \|F_A^-\|^2
  \ge \bigl|\|F_A^+\|^2 - \|F_A^-\|^2\bigr| = 8\pi^2 |c_2(E)|.
\]

\begin{definition_de}[Instanton / Anti-selbst-dualer Zusammenhang]
Ein Zusammenhang heißt \emph{Instanton} (oder \emph{anti-selbst-dual}),
falls $F_A^+ = 0$, d.\,h. $\star F_A = -F_A$.  Instantonen sind absolute
Minima von $\mathcal{YM}$ in ihrer topologischen Klasse.
\end{definition_de}

% =============================================================
\section{Donaldsons Satz und sein Beweis}
% =============================================================

\begin{satz}[Donaldson, 1983 \cite{Donaldson1983}]
Sei $X$ eine glatte, geschlossene, einfach zusammenhängende 4-Mannigfaltigkeit
mit definiter Schnittform $Q_X$.  Dann ist $Q_X$ über $\ZZ$ diagonalisierbar,
d.\,h. $Q_X \cong \pm \mathbf{1}_{b_2}$.
\end{satz}

\begin{proof}[Beweisidee]
Angenommen $Q_X$ ist positiv definit mit $b_2 = b$.  Donaldson betrachtet
den Modulraum $\mathcal{M}_1$ der ASD-Zusammenhänge auf dem Bündel
$E \to X$ mit $c_2(E) = 1$ (Rang 2, $\mathrm{SU}(2)$-Bündel) für eine
generische Riemannsche Metrik $g$.

\medskip
\noindent\textbf{Schritt 1 (Regularität und Dimension).}
Die erwartete reelle Dimension von $\mathcal{M}_1$ ergibt sich aus dem
Atiyah--Singer-Indexsatz für den elliptischen Komplex zur ASD-Deformation:
\[
  \dim \mathcal{M}_1 = 8c_2 - 3(b_0 - b_1 + b_2^+) = 5.
\]
Für eine generische Metrik ist $\mathcal{M}_1$ eine glatte 5-dimensionale
Mannigfaltigkeit abseits reduzibler Zusammenhänge.

\medskip
\noindent\textbf{Schritt 2 (Rand und Kompaktifizierung).}
Nach Uhlenbecks Kompaktheitssatz \cite{Uhlenbeck1982} besteht der ideale
Rand $\partial\overline{\mathcal{M}}_1$ aus Blaseninstantonen an einzelnen
Punkten; topologisch ist dieser Rand homöomorph zu $X$ selbst.

\medskip
\noindent\textbf{Schritt 3 (Kobordismus-Argument).}
Damit ist $\overline{\mathcal{M}}_1$ ein kompakter 5-dimensionaler Kobordismus
zwischen $X$ und einer Vereinigung von $b$ Kopien von $\CP^2$.  Dies erzwingt,
dass $Q_X$ diagonal ist. $\square$
\end{proof}

\begin{korollar}
Keine glatte einfach zusammenhängende 4-Mannigfaltigkeit hat Schnittform $E_8$
oder eine andere nicht-diagonalisierbare definite Form.  Insbesondere besitzt
die \emph{Freedmansche $E_8$-Mannigfaltigkeit} keine glatte Struktur.
\end{korollar}

% =============================================================
\section{Uhlenbecks Kompaktheitssatz}
% =============================================================

\begin{satz}[Uhlenbeck, 1982 \cite{Uhlenbeck1982}]
Sei $\{A_i\}$ eine Folge von ASD-Zusammenhängen auf $E \to X$ mit gleichmäßig
beschränkter Yang--Mills-Energie $\mathcal{YM}(A_i) \le C$.  Dann gibt es:
\begin{enumerate}[label=(\roman*)]
  \item eine endliche Menge von Punkten $\{x_1,\ldots,x_k\} \subset X$
        (den \emph{Blasenpunkten}),
  \item eine Teilfolge,
  \item Eichtransformationen $g_i \in \mathcal{G}|_{X \setminus \{x_j\}}$,
\end{enumerate}
so dass $g_i^* A_i$ in $C^\infty_{\mathrm{lok}}$ auf
$X \setminus \{x_1,\ldots,x_k\}$ gegen einen glatten ASD-Zusammenhang
$A_\infty$ konvergiert.  An jedem Blasenpunkt $x_j$ entsteht in einem
reskalisierten Grenzwert ein Instanton auf $S^4$.
\end{satz}

% =============================================================
\section{Donaldson-Invarianten}
% =============================================================

Für $b_2^+(X) > 1$ definiert Donaldson polynomiale Invarianten
\[
  D_X^k \colon \mathrm{Sym}^{d(k)}(H_2(X;\ZZ)) \longrightarrow \ZZ,
\]
wobei $d(k) = 4k - \tfrac{3}{2}(1 + b_2^+)$,
durch Integration über den Modulraum $\mathcal{M}_k$ von Kohomologieklassen,
die durch die $\mu$-Abbildung $\mu \colon H_2(X;\ZZ) \to H^2(\mathcal{M}_k;\ZZ)$
entstehen.  Die Donaldson-Invarianten sind Diffeomorphieinvarianten von $X$.

% =============================================================
\section{Seiberg--Witten-Theorie}
% =============================================================

1994 führten Seiberg und Witten \cite{SeibergWitten1994} ein weit einfacheres
Gleichungssystem ein, motiviert durch $N{=}2$-supersymmetrische
Yang--Mills-Theorie.

\begin{definition_de}[Seiberg--Witten-Gleichungen]
Sei $X$ glatt, geschlossen und orientiert mit einer $\Spin^c$-Struktur
$\mathfrak{s}$.  Seien $S^\pm \to X$ die Spinorbündel und $L = \det\mathfrak{s}$
das Determinantenlinienbündel.  Die SW-Gleichungen für
$(A,\Phi) \in \mathcal{A}(L) \times \Gamma(S^+)$ sind:
\begin{align}
  D_A \Phi &= 0,\\
  F_A^+ &= \sigma(\Phi),
\end{align}
wobei $D_A$ der durch $A$ getwistete Dirac-Operator und
$\sigma(\Phi) \in i\Omega^+$ die quadratische Abbildung ist.
\end{definition_de}

\begin{definition_de}[Seiberg--Witten-Invariante]
Für $b_2^+(X) > 1$ ist die \emph{Seiberg--Witten-Invariante}
\[
  \SW_X \colon \mathrm{Spin}^c(X) \longrightarrow \ZZ,
\]
definiert durch Zählung (mit Vorzeichen) von Lösungen der SW-Gleichungen
im Modulraum $\mathcal{M}_{SW}(\mathfrak{s})$ der Dimension 0.
\end{definition_de}

\begin{satz}[Grundeigenschaften der SW-Invarianten \cite{Witten1994SW}]
\begin{enumerate}[label=(\roman*)]
  \item $\SW_X(\mathfrak{s}) \neq 0$ nur für endlich viele $\mathfrak{s}$
        (die \emph{Basisklassen}).
  \item Trägt $X$ eine Riemannsche Metrik mit positiver skalarer Krümmung,
        so gilt $\SW_X \equiv 0$.
  \item Für Kähler-Flächen sind die Basisklassen $\pm c_1(K)$.
\end{enumerate}
\end{satz}

\begin{satz}[Wittensche Vermutung, bewiesen von Taubes \cite{Taubes1994}]
Für Mannigfaltigkeiten vom \emph{einfachen Typ} bestimmen die
SW-Invarianten die Donaldson-Invarianten gemäß
\[
  D_X = e^{Q/2} \sum_{\mathfrak{s}:\,\SW(\mathfrak{s})\neq 0} \SW(\mathfrak{s})\,
        e^{c_1(\mathfrak{s})}.
\]
\end{satz}

% =============================================================
\section{Exotische $\RR^4$-Strukturen}
% =============================================================

\begin{satz}[Donaldson--Freedman]
Es gibt überabzählbar viele paarweise nicht-diffeomorphe glatte Strukturen
auf $\RR^4$.  Dies steht im schärfsten Kontrast zu allen anderen Dimensionen
$n \neq 4$, in denen $\RR^n$ eine eindeutige glatte Struktur trägt.
\end{satz}

\begin{bemerkung}
Die Existenz exotischer $\RR^4$ folgt aus der Kombination von Freedmans
topologischem $h$-Kobordismussatz mit Donaldsons Nicht-Diagonalisierbarkeits\-resultat:
Die Freedmansche $E_8$-Mannigfaltigkeit minus einem Punkt enthält ein
exotisches $\RR^4$.
\end{bemerkung}

% =============================================================
\section{Kirby-Kalkül}
% =============================================================

\begin{definition_de}[Henkelzerlegung]
Jede glatte 4-Mannigfaltigkeit $X$ besitzt eine Henkelzerlegung
$X = h^0 \cup h^1_1 \cup \cdots \cup h^2_{k_2} \cup h^3_1 \cup \cdots \cup h^4_1$,
wobei ein $k$-Henkel $h^k \cong D^k \times D^{4-k}$ entlang
$S^{k-1} \times D^{4-k}$ angeklebt wird.
\end{definition_de}

\begin{satz}[Kirby, 1978 \cite{Kirby1978}]
Zwei gerahmte Verschlingungen stellen genau dann diffeomorphe 4-Mannigfaltigkeiten dar,
wenn sie durch eine Folge von Kirby-Zügen (Henkelgleiten und Stabilisierung)
auseinander hervorgehen.
\end{satz}

% =============================================================
\section{Die glatte vierdimensionale Poincaré-Vermutung}
% =============================================================

\begin{vermutung}[Glatte Poincaré-Vermutung in Dimension 4 -- OFFEN]
\label{verm:s4pc}
Jede glatte Homotopie-4-Sphäre ist diffeomorph zu $S^4$.
\end{vermutung}

\begin{bemerkung}
Die topologische Version wurde von Freedman bewiesen.  In allen anderen
Dimensionen folgt die glatte Poincaré-Vermutung aus Chirurgietheorie
(Smale für $n \ge 6$, Perelman für $n = 3$).  Dimension 4 ist der einzige
offen gebliebene Fall.
\end{bemerkung}

Die Hauptkandidaten für potenzielle Gegenbeispiele sind \emph{Gluck-Twists}
$X_K$ auf $S^4$ und gewisse Cappell--Shaneson-Homotopiesphären.
Alle bisher untersuchten Kandidaten erwiesen sich als diffeomorph zu $S^4$
\cite{AkbulutKirby1979}.

% =============================================================
\section{Schnittformen: Einschränkungen aus der Eichtheorie}
% =============================================================

Die Kombination von Donaldsons Satz mit Freedmans Klassifikation liefert:

\begin{korollar}
Eine glatte einfach zusammenhängende 4-Mannigfaltigkeit mit definiter
Schnittform ist homöomorph zu $\#_{b_2} \CP^2$ oder $\#_{b_2} \overline{\CP}^2$.
Ihre Schnittform ist $\pm \mathbf{1}_{b_2}$.
\end{korollar}

\begin{center}
\begin{tabular}{clll}
\toprule
$b_2$ & Form & Topologisch & Glatt \\
\midrule
1 & $(1)$ & $\CP^2$ & $\CP^2$ (eindeutig) \\
1 & $(-1)$ & $\overline{\CP}^2$ & $\overline{\CP}^2$ (eindeutig) \\
8 & $E_8$ & Freedman $E_8$-Mfg. & \emph{keine glatte Struktur} \\
16 & $E_8 \oplus E_8$ & existiert topologisch & \emph{keine glatte Struktur} \\
\bottomrule
\end{tabular}
\end{center}

% =============================================================
\section{Offene Probleme}
% =============================================================

\begin{vermutung}[Glatte Poincaré in Dimension 4, = Vermutung~\ref{verm:s4pc}]
Jede glatte Homotopie-4-Sphäre ist diffeomorph zu $S^4$.
\end{vermutung}

\begin{vermutung}[Minimum-Genus-Problem]
Für eine einfach zusammenhängende 4-Mannigfaltigkeit $X$ wird das minimale
Geschlecht einer glatt eingebetteten Fläche, die eine Klasse
$\alpha \in H_2(X;\ZZ)$ repräsentiert, durch die \emph{Adjunktionsformel}
$g(\alpha) = 1 + \tfrac{1}{2}(\alpha \cdot \alpha + |\alpha \cdot K|)$
gegeben, wobei $K$ die kanonische Klasse ist.  Dies wurde (von Kronheimer--Mrowka
\cite{KronheimerMrowka1994}) für Klassen mit $\alpha \cdot \alpha > 0$
mithilfe der SW-Theorie bewiesen; der allgemeine Fall enthält weitere offene
Fragen.
\end{vermutung}

\begin{vermutung}[Einfacher-Typ-Vermutung]
Jede glatte geschlossene 4-Mannigfaltigkeit mit $b_2^+ > 1$ ist vom
SW-einfachen Typ.
\end{vermutung}

\begin{vermutung}[Geographieproblem]
Welche Paare $(b_2^+, \chi_h)$ werden durch minimale Flächen allgemeinen
Typs realisiert?  Die Bogomolov--Miyaoka--Yau-Ungleichung $c_1^2 \le 3c_2$
und die Noether-Ungleichung liefern notwendige Bedingungen; die vollständige
Geographie bleibt offen.
\end{vermutung}

% =============================================================
\section{Fazit}
% =============================================================

Donaldsons Entdeckung, dass Instantonen topologische Information tragen,
hat sowohl die Mathematik als auch die mathematische Physik transformiert.
Die Nicht-Existenz glatter Strukturen auf gewissen topologischen
4-Mannigfaltigkeiten und die Fülle exotischer glatter Strukturen auf anderen
machen Dimension 4 zu einem Ausnahmefall in der Mannigfaltigkeitstopologie.

Die Seiberg--Witten-Theorie vereinfachte den eichtheoretischen Ansatz,
und Taubes' Identifikation der SW-Invarianten mit
Gromov--Witten-Invarianten für symplektische Mannigfaltigkeiten öffnete
weitere Verbindungen.  Doch die glatte vierdimensionale Poincaré-Vermutung
bleibt hartnäckig offen.

\begin{thebibliography}{99}

\bibitem{Donaldson1983}
S.~K. Donaldson,
\textit{An application of gauge theory to four dimensional topology},
J. Differential Geom. \textbf{18} (1983), Nr.~2, 279--315.

\bibitem{Donaldson1986}
S.~K. Donaldson,
\textit{Connections, cohomology and the intersection forms of 4-manifolds},
J. Differential Geom. \textbf{24} (1986), Nr.~3, 275--341.

\bibitem{DonaldsonKronheimer1990}
S.~K. Donaldson und P.~B. Kronheimer,
\textit{The Geometry of Four-Manifolds},
Oxford Mathematical Monographs, Oxford University Press, 1990.

\bibitem{Freedman1982}
M.~H. Freedman,
\textit{The topology of four-dimensional manifolds},
J. Differential Geom. \textbf{17} (1982), Nr.~3, 357--453.

\bibitem{Uhlenbeck1982}
K.~K. Uhlenbeck,
\textit{Removable singularities in Yang--Mills fields},
Comm. Math. Phys. \textbf{83} (1982), Nr.~1, 11--29.

\bibitem{SeibergWitten1994}
N.~Seiberg und E.~Witten,
\textit{Electric-magnetic duality, monopole condensation, and confinement in
$N=2$ supersymmetric Yang--Mills theory},
Nuclear Phys. B \textbf{426} (1994), 19--52.

\bibitem{Witten1994SW}
E.~Witten,
\textit{Monopoles and four-manifolds},
Math. Res. Lett. \textbf{1} (1994), 769--796.

\bibitem{Taubes1994}
C.~H. Taubes,
\textit{The Seiberg--Witten invariants and symplectic forms},
Math. Res. Lett. \textbf{1} (1994), 809--822.

\bibitem{KronheimerMrowka1994}
P.~B. Kronheimer und T.~S. Mrowka,
\textit{The genus of embedded surfaces in the projective plane},
Math. Res. Lett. \textbf{1} (1994), 797--808.

\bibitem{Kirby1978}
R.~Kirby,
\textit{A calculus for framed links in $S^3$},
Invent. Math. \textbf{45} (1978), 35--56.

\bibitem{AkbulutKirby1979}
S.~Akbulut und R.~Kirby,
\textit{Mazur manifolds},
Michigan Math. J. \textbf{26} (1979), 259--284.

\bibitem{GompfStipsicz1999}
R.~E. Gompf und A.~I. Stipsicz,
\textit{4-Manifolds and Kirby Calculus},
Graduate Studies in Mathematics \textbf{20},
American Mathematical Society, 1999.

\end{thebibliography}

\end{document}
